\section*{Hoofdstuk \ref{ch:b3:submanifolds} --- Deelvariëteiten
van $\R^n$}

\begin{solution}{exo:b3:submanifolds:1}
(a) Elk is $F^{-1}(0)$ voor een submersie: $\norm x^2 - 1$ op
$\R^n\setminus\{0\}$ (gradiënt $2x \neq 0$): dimensie $n-1$; $x^2
+ y^2 - z^2 - 1$ (gradiënt $(2x, 2y, -2z) \neq 0$ op de
hyperboloïde, waar $x^2 + y^2 = 1 + z^2 > 0$): dimensie $2$; de
torusfunctie $G = (\rho - R)^2 + z^2 - r^2$ met $\rho =
\sqrt{x^2+y^2}$ is $\mathcal C^\infty$ nabij de torus (daar is
$\rho \geq R - r > 0$) met $\nabla G \neq 0$ (haar
$z$-component is $2z$, en waar $z = 0$ is de radiale component
$2(\rho - R)\ne0$, want $\abs{\rho - R} = r$): dimensie $2$.

(b) Voor kleine $\varepsilon$ heeft
$(C\setminus\{0\})\cap B(0,\varepsilon)$ precies $2$
\omterm{def:b3:topology:components}{samenhangscomponenten} (de bovenste en de onderste geperforeerde
mantel, elk \omterm{def:b3:topology:pathconnected}{wegsamenhangend}: verbind via cirkels en stralen). Was
$C$ in $0$ een $2$-deelvariëteit, dan zou een rechttrekking een
\omterm{def:b3:topology:continuity}{homeomorfisme} van $C\cap\Omega$ op een open stuk van een vlak
geven dat $0$ naar een punt $p$ stuurt; kleine geperforeerde
vlakomgevingen van $p$ hebben \emph{één} component, en
\omterm{def:b3:topology:continuity}{homeomorfismen} behouden het aantal componenten van geperforeerde
\omterm{def:b3:topology:topology}{omgevingen}: tegenspraak.
\end{solution}

\begin{solution}{exo:b3:submanifolds:2}
(a) $T_aS^{n-1} = \ker\bigl(2a^{\mathsf T}\bigr) = a^\perp$.
(b) In $p = ((R+r)\cos\theta, (R+r)\sin\theta, 0)$ is $\nabla G =
(2r\cos\theta, 2r\sin\theta, 0)$, dus is het raakvlak
$\operatorname{Vect}\bigl((-\sin\theta, \cos\theta, 0),\ (0, 0,
1)\bigr)$: het verticale vlak dat aan de buitenevenaar raakt.
(c) De schroeflijn is een ingebedde kromme met $\varphi'(t_0) =
(-\sin t_0, \cos t_0, 1) \neq 0$: de raaklijn in $\varphi(t_0)$
is $\varphi(t_0) + \R\,\varphi'(t_0)$, zoals
\cref{prop:b3:submanifolds:tangent}(2) voorschrijft.
\end{solution}

\begin{solution}{exo:b3:submanifolds:3}
(a) $Df(x,y) = \bigl(\begin{smallmatrix}2x & -2y\\ 2y &
2x\end{smallmatrix}\bigr)$ met $\det = 4(x^2 + y^2)$: de stelling
is in elk punt behalve de oorsprong van toepassing.
(b) $f(-z) = f(z)$: nooit injectief op een verzameling die
symmetrisch is om $0$; op $\R^2\setminus\{0\}$ is ze overal een
lokaal diffeomorfisme en toch globaal twee-op-één. Nabij $(1, 0)
= f(\pm(1, 0))$ zijn de twee inversen de twee takken van de
vierkantswortel: in complexe notatie $w \mapsto \pm\sqrt w$ (de
hoofdtak), dat wil zeggen
\[
(u, v) \longmapsto \pm\Bigl(\sqrt{\tfrac{u +
\sqrt{u^2+v^2}}{2}},\ \frac{v}{2\sqrt{(u +
\sqrt{u^2+v^2})/2}}\Bigr).
\]
(c) Jacobiaan $r > 0$: een lokaal diffeomorfisme op
$\intoo0\infty\times\R$, maar $\theta \mapsto \theta + 2\pi$
geeft hetzelfde punt: lokaal inverteerbaar (de hoek ligt op een
halfvlak op $2\pi$ na vast), nooit globaal.
\end{solution}

\begin{solution}{exo:b3:submanifolds:4}
(a) $\nabla F = 3(x^2 - y,\ y^2 - x)$ verdwijnt dan en slechts
dan als $y = x^2$ en $x = y^2$, dat wil zeggen $x^4 = x$: in
$(0,0)$ en $(1,1)$; enkel $(0,0)$ ligt op $\mathcal C$ (want
$F(1,1) = -1$). Op $\mathcal C\setminus\{0\}$ is $F$ dus een
submersie: een $1$-deelvariëteit, lokaal een grafiek $y =
\psi(x)$ waar $F_y = 3(y^2 - x) \neq 0$ en $x = \chi(y)$ waar
$F_x = 3(x^2 - y) \neq 0$ (buiten de oorsprong geldt er minstens
één van beide).

(b) In $(\frac32, \frac32)$: $\nabla F = 3(\frac94 - \frac32)(1,
1) = \frac94(1,1)$: raaklijn $x + y = 3$.

(c) De rationale parametrisering $x = \frac{3t}{1 + t^3}$, $y =
\frac{3t^2}{1+t^3}$ gaat in $t = 0$ door $0$ met snelheid $(3,
0)$; $x \leftrightarrow y$ verwisselen (de kromme is symmetrisch,
of herparametriseer met $1/t$) geeft een tweede $\mathcal
C^1$-kromme door $0$ met snelheid $(0, 3)$. Twee onafhankelijke
raakrichtingen zijn onmogelijk voor een $1$-deelvariëteit (haar
\omterm{def:b3:submanifolds:tangent}{raakruimte} is een rechte,
\cref{prop:b3:submanifolds:tangent}): $\mathcal C$ is in de
oorsprong geen deelvariëteit --- een transversale
zelfkruising.
\end{solution}

\begin{solution}{exo:b3:submanifolds:5}
(a) $F(M + H) = M^{\mathsf T}M + M^{\mathsf T}H + H^{\mathsf
T}M + H^{\mathsf T}H$: dus $DF(M)(H) = M^{\mathsf T}H +
H^{\mathsf T}M$. Voor $M \in O_n$ en symmetrische $S$ geeft $H =
\frac12MS$ dat $DF(M)(H) = \frac12(S + S^{\mathsf T}) = S$:
surjectief op $S_n$.

(b) $O_n = F^{-1}(I)$ met $F$ in elk van haar punten een
submersie (op $S_n$, van dimensie $\frac{n(n+1)}2$): dus een
deelvariëteit van dimensie $n^2 - \frac{n(n+1)}2 =
\frac{n(n-1)}2$. \omterm{def:b3:topology:compact}{Compact}: gesloten ($F$ is \omterm{def:b3:topology:continuity}{continu}) en begrensd
(de kolommen zijn eenheidsvectoren). \omterm{def:b3:submanifolds:tangent}{Raakruimte} in $I$: $\ker
DF(I) = \{H : H + H^{\mathsf T} = 0\}$.

(c) $\bigl(\eu^{tH}\bigr)^{\mathsf T}\eu^{tH} =
\eu^{tH^{\mathsf T}}\eu^{tH} = \eu^{-tH}\eu^{tH} = I$ (het
transponeren gaat door de reeks heen; de exponentiëlen van
commuterende matrices vermenigvuldigen,
\cref{thm:b3:ode:matrixexp}).
\end{solution}

\begin{solution}{exo:b3:submanifolds:6}
(a) $\det(M + H) = \det M\,\det(I + M^{-1}H) = \det M\bigl(1 +
\operatorname{tr}(M^{-1}H) + O(\norm H^2)\bigr)$ voor
inverteerbare $M$ (de ontwikkeling van $\det$ nabij $I$: de
lineaire term van $\prod(1 + \lambda_i)$); en met $\det M\cdot
M^{-1} = \operatorname{com}(M)^{\mathsf T}$: $D\det(M)(H) =
\operatorname{tr}\bigl(\operatorname{com}(M)^{\mathsf T}H\bigr)$,
en de formule breidt zich met \omterm{ex:b3:lebesgue:gamma}{dichtheid} en \omterm{def:b3:topology:continuity}{continuïteit} uit tot
alle $M$. Op $SL_n$ is $\det M = 1$: $D\det(M) \ne 0$ (haar
waarde op $H = M$ is $\operatorname{tr}(I)\cdot\dots = n\det M =
n$).

(b) $SL_n = \det^{-1}(1)$ met $\det$ daar een submersie (met
waarden in $\R$): dimensie $n^2 - 1$; en $T_ISL_n = \ker
D\det(I) = \{H : \operatorname{tr}H = 0\}$.

(c) $GL_n$ is een \emph{open} deelverzameling van $M_n(\R)$
(\cref{exo:b3:topology:8}): een deelvariëteit van volle dimensie
$n^2$ (rechttrekking: de identieke kaart).
\end{solution}

\begin{solution}{exo:b3:submanifolds:7}
(a) $(y, x) = \lambda(2x, 2y)$ en $x^2 + y^2 = 1$: uit $y =
2\lambda x$ en $x = 2\lambda y$ volgt $x^2 = y^2 = \frac12$.
Waarden van $xy$: $\pm\frac12$: maximum $\frac12$ in
$\pm\frac1{\sqrt2}(1,1)$, minimum $-\frac12$ in
$\pm\frac1{\sqrt2}(1,-1)$ (de nevenvoorwaarde geeft een \omterm{def:b3:topology:compact}{compacte}
verzameling: de extrema bestaan).

(b) Op het inwendige van de simplex ($p_i > 0$) geeft Lagrange
voor $H(p) = -\sum p_i\ln p_i$ met nevenvoorwaarde $\sum p_i = 1$
dat $-\ln p_i - 1 = \lambda$ voor alle $i$: alle $p_i$ gelijk,
dus $p_i = \frac1n$, met $H = \ln n$. Het maximum over de
\omterm{def:b3:topology:compact}{compacte} simplex wordt bereikt; werd het op de rand bereikt (een
$p_i = 0$), dan leeft de verdeling op $\leq n - 1$ punten en is
met inductie $H \leq \ln(n-1) < \ln n$: het inwendige kritieke
punt is dus het globale maximum --- uniforme onwetendheid
maximaliseert de entropie.

(c) Minimaliseer $(x-1)^2 + y^2$ op de \omterm{def:b3:topology:compact}{compacte} ellips:
$(2(x{-}1), 2y) = \lambda(\frac x2, 2y)$. Is $y \neq 0$, dan is
$\lambda = 1$, en geeft $2(x - 1) = \frac x2$ dat $x = \frac43$
en $y^2 = 1 - \frac49 = \frac59$: afstand$^2$ $= \frac19 +
\frac59 = \frac23$. Is $y = 0$, dan is $x = \pm2$, met afstanden
$1$ en $3$. Dichtstbijzijnde punten: $\bigl(\frac43,
\pm\frac{\sqrt5}3\bigr)$, op afstand $\sqrt{2/3} < 1$. De
vergelijking met de multiplicator zegt dat het lijnstuk van
$(1,0)$ naar het dichtstbijzijnde punt evenwijdig is met
$\nabla$(ellips): het snijdt de ellips loodrecht, zoals de
meetkunde eist.
\end{solution}

\begin{solution}{exo:b3:submanifolds:8}
Inductie naar $n$; $n = 1$ is triviaal. De functie van Rayleigh
$f(x) = \langle Ax, x\rangle$ bereikt haar maximum $\lambda_1$ op
de \omterm{def:b3:topology:compact}{compacte} $S^{n-1}$ in een $v_1$; Lagrange
(\cref{thm:b3:submanifolds:lagrange}, de sfeer als
niveauverzameling) geeft $2Av_1 = 2\lambda v_1$, en $\lambda =
\langle Av_1, v_1\rangle = \lambda_1$. Het hypervlak $v_1^\perp$
is $A$-invariant (symmetrie: $\langle Av, v_1\rangle = \langle v,
Av_1\rangle = 0$); de beperking is symmetrisch, en inductie
levert een orthonormale eigenbasis $v_2, \dots, v_n$ van
$v_1^\perp$ met eigenwaarden $\lambda_2 \geq \dots \geq
\lambda_n$, elk het maximum van $f$ op de sfeer van het
resterende \omterm{thm:b3:hilbert:decomposition}{orthogonaal complement}. Courant--Fischer volgt precies
als in \cref{exo:b3:spectral:8}: ontwikkel $x = \sum c_iv_i$;
snijd op een $k$-dimensionale testruimte met
$\operatorname{Vect}(v_k, \dots, v_n)$ (het tellen van dimensies
in $\R^n$) om $\min \leq \lambda_k$ te krijgen, en
$\operatorname{Vect}(v_1, \dots, v_k)$ verwezenlijkt $\min =
\lambda_k$.
\end{solution}

\begin{solution}{exo:b3:submanifolds:9}
Elke kolom tot eenheidsnorm schalen deelt $\abs{\det}$ door
$\prod\norm{c_i}$: het volstaat dus $\abs{\det M} \leq 1$ te
bewijzen als alle kolommen eenheidsvectoren zijn, met gelijkheid
dan en slechts dan als $M \in O_n$. De functie $\det$ is \omterm{def:b3:topology:continuity}{continu}
op de \omterm{def:b3:topology:compact}{compacte} $(S^{n-1})^n$: ze bereikt een maximum $m \geq \det
I = 1 > 0$ in een $M$. Houd alle kolommen op de $i$-de na vast:
dan is $\det$ lineair in $c_i$ met als gradiënt de $i$-de kolom
van $\operatorname{com}(M)$; Lagrange op de $i$-de sfeer geeft
$\operatorname{com}(M)_{\cdot i} = \lambda_i c_i$. De identiteit
$M^{\mathsf T}\operatorname{com}(M) = \det(M)\,I$ luidt $\langle
c_j, \operatorname{com}(M)_{\cdot i}\rangle = \det(M)\,
\delta_{ij}$, dat wil zeggen $\lambda_i\langle c_j, c_i\rangle =
\det(M)\delta_{ij}$; met $j = i$: $\lambda_i = \det M = m \neq
0$, en dan geeft $j \neq i$ dat $\langle c_i, c_j\rangle = 0$: de
kolommen zijn orthonormaal, dus $M \in O_n$ en $m = \abs{\det M}
= 1$. Bijgevolg is altijd $\abs{\det} \leq \prod\norm{c_i}$, met
gelijkheid precies voor orthogonale kolommen (schaal terug): het
volume van een parallellepipedum is bij gegeven ribbelengten het
grootst wanneer de ribben loodrecht staan.
\end{solution}

\begin{solution}{exo:b3:submanifolds:10}
$y = \pm\sqrt{1 - x^2}$ voldoet nabij elk punt met $y \neq 0$; $x
= \pm\sqrt{1 - y^2}$ nabij elk punt met $x \neq 0$; in $(1, 0)$:
de grafiek $x = \sqrt{1 - y^2}$ boven $y \in \intoo{-1}1$, wat
\cref{thm:b3:submanifolds:characterizations}(3) is met de
coördinaten verwisseld. Een of andere permutatie werkt altijd
omdat de raaklijn, eendimensionaal als ze is, niet tegelijk
verticaal en horizontaal kan zijn --- maar ze kan wel elk van
beide zijn, zodat geen vaste keuze van ``afhankelijke''
coördinaat in elk punt dienst doet.
\end{solution}

\begin{solution}{exo:b3:submanifolds:11}
(a) De definiërende afbeelding $F(M) = M^{\mathsf T}M - I$ is
\omterm{def:b3:topology:continuity}{continu}: $O_n = F^{-1}(0)$ is gesloten; de kolommen van een
orthogonale matrix zijn eenheidsvectoren, dus is de euclidische
(frobenius)norm precies $\sqrt n$: begrensd. \omterm{def:b3:topology:compact}{Compact} volgens
Heine--Borel in $M_n(\R) \cong \R^{n^2}$.

(b) $O_n$ is de niveauverzameling $F = 0$ uit het hoofdstuk:
$DF(A)H = A^{\mathsf T}H + H^{\mathsf T}A$, in elke $A \in O_n$
surjectief op de symmetrische matrices (neem bij symmetrische $S$
$H = \frac12AS$), dus is $O_n$ een deelvariëteit van dimensie
$n^2 - \frac{n(n+1)}2 = \frac{n(n-1)}2$ met
\[
T_AO_n = \ker DF(A) = \{H : A^{\mathsf T}H \text{
antisymmetrisch}\} = \{AK : K^{\mathsf T} = -K\} ;
\]
in $A = I$ zijn dat de antisymmetrische matrices.

(c) $\exp(tK)^{\mathsf T}\exp(tK) = \exp(tK^{\mathsf T})
\exp(tK) = \exp(-tK)\exp(tK) = I$ (de twee matrices $\pm tK$
commuteren, dus is het product van de exponentiëlen de
exponentiële van de som): dus $\exp(tK) \in O_n$, en $\gamma(t) =
A\exp(tK)$ is een kromme in $O_n$ met $\gamma(0) = A$ en
$\gamma'(0) = AK$. Terwijl $K$ de antisymmetrische matrices
doorloopt, veegt $AK$ $T_AO_n$ uit: de exponentiële verwezenlijkt
de hele \omterm{def:b3:submanifolds:tangent}{raakruimte} met expliciete krommen --- de sluiproute van
de liegroepen die \cref{pb:b3:submanifolds:1} voor $SO(3)$
uitbuit.
\end{solution}

\begin{solution}{exo:b3:submanifolds:12}
Bestaan: snijd $M$ met een grote gesloten bal rond $p$ om een
niet-lege \omterm{def:b3:topology:compact}{compacte} verzameling te krijgen; de \omterm{def:b3:topology:continuity}{continue} afstand
bereikt daar haar minimum, en punten buiten de bal liggen verder.
Voorwaarde van eerste orde: voor een kromme $\gamma$ in $M$ met
$\gamma(0) = x_0$ is de functie $h(t) = \norm{\gamma(t) - p}^2$
differentieerbaar met een minimum in $0$:
\[
0 = h'(0) = 2\,\langle\gamma'(0),\ x_0 - p\rangle,
\]
en $\gamma'(0)$ veegt $T_{x_0}M$ uit: dus $p - x_0 \perp
T_{x_0}M$. Sfeer $S(c, r)$: de \omterm{def:b3:submanifolds:tangent}{raakruimte} in $x_0$ is $(x_0 -
c)^\perp$, dus $p - x_0 \parallel x_0 - c$: $x_0$ ligt op de
rechte door $c$ en $p$, op afstand $r$ van $c$ --- het punt op de
straal, zoals de meetkunde eist. Parabool: in $x_0 = (x, x^2)$
wordt de raaklijn opgespannen door $(1, 2x)$; de orthogonaliteit
op $p - x_0 = (2 - x, -x^2)$ luidt
\[
(2 - x) - 2x^3 = 0,
\qquad\text{dat wil zeggen}\qquad 2x^3 + x - 2 = 0,
\]
met één reële wortel ($x \mapsto 2x^3 + x$ is strikt stijgend)
$x \approx 0.835$; dan is $x_0 \approx (0.835, 0.698)$ en $d(p,
M) = \sqrt{(2 - 0.835)^2 + 0.698^2} \approx 1.358$.
\end{solution}

\begin{solution}{pb:b3:submanifolds:1}
\textbf{1.} Beeld $q = t + x\mathrm i + y\mathrm j + z\mathrm k
\mapsto \bigl(\begin{smallmatrix}\alpha & \beta\\ -\bar\beta
& \bar\alpha\end{smallmatrix}\bigr)$ af met $\alpha = t + \iu x$
en $\beta = y + \iu z$: je gaat na dat $1, \mathrm i, \mathrm j,
\mathrm k$ naar $I$, $\bigl(\begin{smallmatrix}
\iu & 0\\ 0 & -\iu\end{smallmatrix}\bigr)$,
$\bigl(\begin{smallmatrix}0 & 1\\ -1 &
0\end{smallmatrix}\bigr)$ en $\bigl(\begin{smallmatrix}0 &
\iu\\ \iu & 0\end{smallmatrix}\bigr)$ gaan, waarvan de producten
de quaternionentabel weergeven: de afbeelding is een injectief
algebramorfisme, dus erft $\mathbb H$ de associativiteit; $N(q) =
\abs\alpha^2 + \abs\beta^2 = \det$ is multiplicatief, en de
toevoeging stemt overeen met de getransponeerde adjunct, wat
$\overline{pq} = \bar q\bar p$ geeft. \omterm{ex:b3:groups:actions}{Centrum}: commuteren met
$\mathrm i$ dwingt $y = z = 0$ af, met $\mathrm j$ dwingt $x = 0$
af: dus $\R$.

\textbf{2.} $q\bar q = N(q)$: voor $q \neq 0$ is $q^{-1} = \bar
q/N(q)$: een delingsalgebra. Op $S^3$: $N(pq) = 1$ en
$N(q^{-1}) = 1$: een groep; en $S^3 \subseteq \R^4$ is de
eenheidssfeer: een \omterm{def:b3:topology:compact}{compacte} $3$-deelvariëteit.

\textbf{3.} $v$ is puur dan en slechts dan als $\bar v = -v$; dan
is $\overline{qv\bar q} = q\bar v\bar q = -qv\bar q$: dus behoudt
$\rho_q$ $P$. De lineariteit is duidelijk; en $N(qv\bar q) =
N(q)N(v)N(q) = N(v)$: een isometrie van $(P, N) \cong (\R^3,
\norm\cdot^2)$: dus $\rho_q \in O(3)$. En $\rho_{pq}(v) =
pqv\overline{pq} = p(qv\bar q)\bar p = \rho_p(\rho_q(v))$: een
morfisme.

\textbf{4.} $\rho_q = \mathrm{id}$ dan en slechts dan als $qv =
vq$ voor alle pure $v$, dan en slechts dan als $q$ met $\mathrm
i, \mathrm j, \mathrm k$ commuteert, dat wil zeggen als $q$
centraal is (vraag 1): $q \in \R\cap S^3 = \{\pm1\}$.

\textbf{5.} Schrijf $q = t + p$ ($t \in \R$, $p$ puur): $1 = N(q)
= t^2 + N(p)$, dus $t = \cos\frac\theta2$ en $p =
\sin\frac\theta2\,u$ met $N(u) = 1$ voor een $\theta$ (is $p =
0$, dan werkt $q = \pm1$ triviaal). Omdat $u^2 = -N(u) = -1$,
commuteren $q$ en $u$, en is $\rho_q(u) = qu\bar q = uq\bar q =
u$: de as. Voor pure eenheden $w \perp u$ geeft de productregel
$vw = -\langle v, w\rangle + v\wedge w$ (ontwikkel in coördinaten
met de tabel) dat $uw = u\wedge w = -wu$. Dan is
\[
\rho_q(w) = \bigl(\cos\tfrac\theta2 +
\sin\tfrac\theta2u\bigr)\,w\,\bigl(\cos\tfrac\theta2 -
\sin\tfrac\theta2u\bigr)
= \cos\theta\,w + \sin\theta\,(u\wedge w),
\]
met $uwu = -u^2w = w$ en de formules voor de dubbele hoek: de
rotatie over hoek $\theta$ in het georiënteerde vlak $(w, u\wedge
w)$.

\textbf{6.} Elke $\rho_q$ is een rotatie om $u$ over $\theta$: in
de orthonormale basis $(u, w, u\wedge w)$ heeft haar matrix
determinant $+1$: dus $\operatorname{im}\rho \subseteq SO(3)$.
Surjectiviteit: een matrix $R \in SO(3)$ heeft $1$ als
eigenwaarde, want
\[
\det(R - I) = \det R\,\det(I - R^{\mathsf T}) = \det(I - R)
= (-1)^3\det(R - I),
\]
dus $\det(R - I) = 0$. Neem een eigenvector $u$ met norm $1$; $R$
behoudt $u^\perp$ en beperkt zich daar tot een rotatie over een
hoek $\theta$ (vlak orthogonaal, determinant $1$): dus $R =
\rho_q$ met $q = \cos\frac\theta2 + \sin\frac\theta2\,u$. Met
vraag 4 en de eerste isomorfiestelling
(\cref{thm:b3:groups:firstiso}): $SO(3) \cong S^3/\{\pm1\}$.

\textbf{7.} $O_3$ is een \omterm{def:b3:topology:compact}{compacte} $3$-dimensionale deelvariëteit
(\cref{exo:b3:submanifolds:5}); $\det$ is erop \omterm{def:b3:topology:continuity}{continu} met
waarden in $\{\pm1\}$, dus is $SO(3) = O_3\cap\{\det = 1\}$ open
en gesloten in $O_3$: een vereniging van \omterm{def:b3:topology:components}{samenhangscomponenten},
dus zelf een \omterm{def:b3:topology:compact}{compacte} $3$-deelvariëteit, met dezelfde \omterm{def:b3:submanifolds:tangent}{raakruimte}
in $I$: de antisymmetrische matrices.

\textbf{8.} $A_u^2v = u\wedge(u\wedge v) = \langle u, v\rangle u
- v$ (met $u$ van norm $1$), dus $A_u^3v = u\wedge(\langle
u,v\rangle u - v) = -u\wedge v$: $A_u^3 = -A_u$. De exponentiële
reeks splitsen naar de machten modulo de betrekking $A^3 = -A$:
\[
\eu^{\theta A_u} = I + \Bigl(\theta - \frac{\theta^3}{3!} +
\cdots\Bigr)A_u + \Bigl(\frac{\theta^2}{2!} -
\frac{\theta^4}{4!} + \cdots\Bigr)A_u^2
= I + \sin\theta\,A_u + (1 - \cos\theta)\,A_u^2 .
\]
Op $u$: $A_uu = 0$: vastgelaten. Op $w \perp u$: $\eu^{\theta
A_u}w = w + \sin\theta\,u\wedge w + (1 - \cos\theta)(-w) =
\cos\theta\,w + \sin\theta\,u\wedge w$: de rotatie met as $u$
over hoek $\theta$ --- Rodrigues. Elke rotatie heeft die vorm
(vraag 6): $\exp$ is surjectief van de antisymmetrische matrices
op $SO(3)$.

\textbf{9.} Met $q(t) = \cos\frac t2 + \sin\frac t2\,u$ toont
vraag 5 dat $\rho_{q(t)}$ de rotatie met as $u$ over hoek $t$ is,
dat wil zeggen $\rho_{q(t)} = \eu^{tA_u}$, waarvan de afgeleide
in $t = 0$ gelijk is aan $A_u$. Het quaternion loopt op de halve
hoek --- het analytische spoor van de dubbele overdekking.

\textbf{10.} $\rho_{q(t)}$ is de rotatie om $\mathrm k$ over hoek
$t$: in $t = 2\pi$ keert ze terug naar de identiteit --- een lus
in $SO(3)$. Haar lift voldoet aan $q(2\pi) = \cos\pi = -1 =
-q(0)$: de gelifte weg is \emph{niet} gesloten; pas in $t = 4\pi$
keert $q$ terug naar $1$. Interpretatie: de lus van volledige
rotaties is in $SO(3)$ niet samentrekbaar --- haar lift eindigt
op het andere blad van de overdekking --- terwijl de dubbele lus
dat wel is; een lichaam dat met riemen aan zijn \omterm{def:b3:topology:topology}{omgeving} vastzit
(de riemtruc) keert na $4\pi$ terug naar een onverdraaide
toestand, maar niet na $2\pi$. Rotatiegroepen onthouden de
pariteit van hele omwentelingen; en $S^3$, dat \omterm{def:b3:conformal:simplyconnected}{enkelvoudig
samenhangend} is, is waar dat geheugen woont.

\textbf{11.} Kwartslagen: $q_{\mathrm i} = \cos\frac\pi4 +
\sin\frac\pi4\,\mathrm i$ en $q_{\mathrm j} = \cos\frac\pi4 +
\sin\frac\pi4\,\mathrm j$. Product (met de $\mathrm j$-slag
eerst):
\[
q_{\mathrm i}q_{\mathrm j} = \tfrac12(1 + \mathrm i)(1 +
\mathrm j) = \tfrac12\bigl(1 + \mathrm i + \mathrm j +
\mathrm k\bigr),
\]
met norm $1$ en $\cos\frac\theta2 = \frac12$: dus $\theta =
\frac{2\pi}3$, met as $u = \frac{\mathrm i + \mathrm j + \mathrm
k}{\sqrt3}$ (het genormeerde pure deel). Twee opeenvolgende
kwartslagen om orthogonale assen vormen een rotatie over
$120^\circ$ om de hoofddiagonaal van de kubus --- vier reële
vermenigvuldigingen aan boekhouding, met de orthogonaliteit exact
behouden: daarom stellen vluchtsoftware en grafische motoren
rotaties met quaternionen samen.

\textbf{12.} Uit de tabel is $\mathrm{ji} = -\mathrm k$ en
$\mathrm{ki} = \mathrm j$, dus
\[
q\,\mathrm i = a\mathrm i - b + c(\mathrm{ji}) +
d(\mathrm{ki}) = -b + a\mathrm i + d\mathrm j - c\mathrm k .
\]
Vermenigvuldigen met $\bar q = a - b\mathrm i - c\mathrm j -
d\mathrm k$ met de regel voor scalair en vector $(t_1 + p_1)(t_2
+ p_2) = t_1t_2 - \langle p_1, p_2\rangle + t_1p_2 + t_2p_1 +
p_1\wedge p_2$, met $p_1 = (a, d, -c)$ en $p_2 = (-b, -c, -d)$:
het scalaire deel is $-ab - (-ab) = 0$ (puur, zoals het hoort),
en het vectordeel is
\[
b(b, c, d) + a(a, d, -c) + (-c^2 - d^2,\ bc + ad,\ bd - ac)
= \bigl(a^2 + b^2 - c^2 - d^2,\ 2(bc + ad),\ 2(bd -
ac)\bigr):
\]
de eerste kolom van $R_q$. De cyclische afbeelding
$\sigma(\mathrm i) = \mathrm j$, $\sigma(\mathrm j) = \mathrm k$,
$\sigma(\mathrm k) = \mathrm i$ behoudt de betrekkingen $\mathrm
i^2 = \mathrm j^2 = \mathrm k^2 = \mathrm{ijk} = -1$ (het woord
$\mathrm{ijk}$ is cyclisch invariant op de betrekking
$\mathrm{ijk} = \mathrm{jki}$ na, die in elke ring geldt: vervoeg
$\mathrm{ijk} = -1$ met de inverteerbare $\mathrm i$), dus breidt
$\sigma$ uit tot een $\R$-algebra-automorfisme, en is
$\sigma(\rho_q(v)) = \rho_{\sigma(q)}(\sigma(v))$. Uitgewerkt is
het beeld van $\mathrm j$ de formule voor de eerste kolom na de
substitutie $(b, c, d) \to (c, d, b)$ met de basis herbenoemd
$\mathrm i \to \mathrm j \to \mathrm k \to \mathrm i$, wat
precies de weergegeven tweede kolom is; nog een slag geeft de
derde.

\textbf{13.} De diagonaal optellen geeft $\operatorname{tr}R_q =
3a^2 - (b^2 + c^2 + d^2) = 4a^2 - 1$ (eenheidsnorm), en met $a =
\cos\frac\theta2$: $4\cos^2\frac\theta2 - 1 = 1 + 2\cos\theta$.
Antisymmetrisch deel: de drie onafhankelijke elementen van $R_q -
R_q^{\mathsf T}$ zijn $4ab, 4ac, 4ad$ (op de posities $(3,2),
(1,3), (2,1)$), dus $\frac12(R_q - R_q^{\mathsf T}) = A_m$ met $m
= 2a\,(b, c, d) = 2\cos\frac\theta2\sin\frac\theta2\,u =
\sin\theta\,u$. Algoritme: $\theta = \arccos
\frac{\operatorname{tr}R - 1}{2} \in \intcc0\pi$; is $0 < \theta
< \pi$, lees dan $u$ af van $\frac{R - R^{\mathsf
T}}{2\sin\theta}$ en stel $q = \pm(\cos\frac\theta2 +
\sin\frac\theta2 u)$; is $\theta = 0$, dan $q = \pm1$. Voor
$\theta = \pi$: $a = 0$, en Rodrigues (vraag 8) geeft $R = I +
2A_u^2 = I + 2(uu^{\mathsf T} - I) = 2uu^{\mathsf T} - I$, dat
wil zeggen $R + I = 2uu^{\mathsf T}$; elke kolom $\neq 0$ van $R
+ I$ is, genormeerd, $\pm u$, en $q = \pm u$.

\textbf{14.} Met $a = b = c = d = \frac12$ verdwijnen alle
diagonaalelementen, en is $2(bc - ad) = 0$, $2(bd + ac) = 1$,
$2(bc + ad) = 1$, $2(cd - ab) = 0$, $2(bd - ac) = 0$, $2(cd + ab)
= 1$:
\[
R_q = \begin{pmatrix} 0 & 0 & 1\\ 1 & 0 & 0\\ 0 & 1 & 0
\end{pmatrix},
\]
de cyclische permutatie $e_1 \to e_2 \to e_3 \to e_1$. Haar spoor
is $0 = 1 + 2\cos\theta$, dus $\theta = \frac{2\pi}3$, en ze laat
$(1,1,1)$ vast: de rotatie over $120^\circ$ om de hoofddiagonaal
--- precies het antwoord van vraag 11, nu zichtbaar als de matrix
die de coördinaatassen cyclisch verwisselt.

\textbf{15.} Op de basis ga je $\Phi(\bar q) = \Phi(q)^*$ na (de
matrix van $\bar q$ heeft $\alpha' = \bar\alpha$ en $\beta' =
-\beta$, wat de toegevoegd getransponeerde is van
$\bigl(\begin{smallmatrix}\alpha & \beta\\ -\bar\beta &
\bar\alpha\end{smallmatrix}\bigr)$). Bijgevolg is
$\Phi(q)^*\Phi(q) = \Phi(\bar qq) = N(q)I$ en $\det\Phi(q) =
\abs\alpha^2 + \abs\beta^2 = N(q)$: voor $q \in S^3$ is dus
$\Phi(q) \in SU(2)$, en $\Phi$ is een injectief morfisme (vraag
1). Surjectiviteit: zij $U = \bigl(\begin{smallmatrix}\alpha &
\beta\\ \gamma & \delta\end{smallmatrix}\bigr)$ met $\det U = 1$;
dan is $U^{-1} = \bigl(\begin{smallmatrix}\delta & -\beta\\
-\gamma & \alpha\end{smallmatrix}\bigr)$, en $U^{-1} = U^* =
\bigl(\begin{smallmatrix}\bar\alpha & \bar\gamma\\ \bar\beta &
\bar\delta\end{smallmatrix}\bigr)$ dwingt $\delta = \bar\alpha$
en $\gamma = -\bar\beta$ af, en dan is $1 = \det U =
\abs\alpha^2 + \abs\beta^2$: dus $U = \Phi(q)$ voor het
eenheidsquaternion $q$ met coördinaten $\alpha = a + \iu b$ en
$\beta = c + \iu d$. Dus $S^3 \cong SU(2)$.

\textbf{16.} Reële scalairen zijn centraal en $N(p) = 1$ geeft
$\overline{pq\bar p} = p\bar q\bar p$, dus $pq\bar p +
\overline{pq\bar p} = p(q + \bar q)\bar p = q + \bar q$: het
reële deel is invariant. Omgekeerd, zij $\operatorname{Re}q =
\operatorname{Re}q' = a$; dan hebben de pure delen dezelfde norm
$\sqrt{1 - a^2} = s$. Is $s = 0$, dan is $q = q' = \pm1$. Is $s >
0$, schrijf dan $q = a + su$ en $q' = a + su'$ met $u, u'$ pure
eenheden; vraag 6 levert een rotatie die $u'$ naar $u$ brengt,
dat wil zeggen een $p \in S^3$ met $\rho_p(u') = u$, en dan is
$pq'\bar p = a + s\rho_p(u') = q$. De klassen van $S^3$ zijn dus
$\{1\}$, $\{-1\}$, en voor elke $a \in \intoo{-1}1$ de $2$-sfeer
$\{a + su : u \text{ pure eenheid}\}$ met straal $s$. Onder
$\Phi$ is $\operatorname{tr}\Phi(q) = \alpha + \bar\alpha =
2\operatorname{Re}q$: de klassen van $SU(2)$ zijn de
niveauverzamelingen van het spoor. Projecteren: is $q' = pq\bar
p$, dan is $\rho_{q'} = \rho_p\rho_q\rho_p^{-1}$; omgekeerd
dwingt $\rho_{q'} = \rho_p\rho_q\rho_p^{-1} = \rho_{pq\bar p}$
dat $q' = \pm pq\bar p$ (de kern), dus $\operatorname{Re}q' =
\pm\operatorname{Re}q$, dat wil zeggen $\cos\frac{\theta'}2 =
\abs{\operatorname{Re}q'} = \abs{\operatorname{Re}q} =
\cos\frac\theta2$ voor de hoeken in $\intcc0\pi$: geconjugeerde
rotaties hebben gelijke hoeken. Omgekeerd laten gelijke hoeken
vertegenwoordigers met hetzelfde niet-negatieve reële deel toe,
geconjugeerd volgens het bovenstaande: in $SO(3)$ is de
conjugatieklasse van een rotatie precies haar hoek.

\textbf{17.} $N$ is multiplicatief, dus is $\abs\cdot$ een
multiplicatieve norm op $\mathbb H \cong \R^4$ en is $\abs{q^k} =
\abs q^k$: de reeks $\sum q^k/k!$ convergeert absoluut in de
eindigdimensionale (en dus volledige) ruimte, gedomineerd door
$\sum\abs q^k/k! = \eu^{\abs q}$. Voor een pure eenheid $u$ is
$u^2 = -1$, dus $(\theta u)^{2m} = (-1)^m\theta^{2m}$ en $(\theta
u)^{2m+1} = (-1)^m\theta^{2m+1}u$; de reeks splitsen geeft
\[
\exp(\theta u) = \sum_m\frac{(-1)^m\theta^{2m}}{(2m)!} +
u\sum_m\frac{(-1)^m\theta^{2m+1}}{(2m+1)!} = \cos\theta +
\sin\theta\,u .
\]
Elke $q \in S^3$ is $\cos\alpha + \sin\alpha\,u$ met $\alpha \in
\intcc0\pi$ (vraag 5): dus $q = \exp(\alpha u)$ en $\exp(P) =
S^3$. Ten slotte is $\exp(su) = \cos s + \sin s\,u = q(2s)$ in de
notatie van vraag 9, en daar was $\rho_{q(t)} = \eu^{tA_u}$: dus
$\rho_{\exp(su)} = \eu^{2sA_u}$.

\textbf{18.} $vw$ coördinaatsgewijs met de tabel ontwikkelen: de
producten $\mathrm i\cdot\mathrm i = -1$, \dots\ geven de
scalair $-(v_1w_1 + v_2w_2 + v_3w_3)$, en de gemengde producten
($\mathrm{ij} = \mathrm k$, $\mathrm{ji} = -\mathrm k$, \dots)
geven de vector $(v_2w_3 - v_3w_2,\ v_3w_1 - v_1w_3,\ v_1w_2 -
v_2w_1)$: dus $vw = -\langle v, w\rangle + v\wedge w$. Het
omgekeerde product aftrekken: $vw - wv = 2\,v\wedge w$ (de
scalaire delen heffen elkaar op, de uitwendige producten tellen
op). Voor de matrices, met $a\wedge(b\wedge c) = b\langle a,
c\rangle - c\langle a, b\rangle$:
\[
[A_v, A_w]x = v\wedge(w\wedge x) - w\wedge(v\wedge x)
= w\langle v, x\rangle - v\langle w, x\rangle
= (v\wedge w)\wedge x = A_{v\wedge w}x .
\]
Afgeleide: $\overline{\exp(tv)} = \exp(-tv)$ (de toevoeging is
\omterm{def:b3:topology:continuity}{continu} en keert het teken van pure quaternionen om), dus
\[
\frac{\dd}{\dd t}\Bigr|_{t=0}\exp(tv)\,x\,\exp(-tv) = vx -
xv = 2\,v\wedge x = 2A_vx :
\]
de differentiaal is $v \mapsto 2A_v$, een lineaire bijectie van
$P$ op de antisymmetrische matrices, en $[2A_v, 2A_w] =
4A_{v\wedge w} = 2A_{2v\wedge w} = 2A_{[v,w]}$ toont dat ze de
\omterm{def:b3:groups:derived}{commutator} van de quaternionen omzet in die van de matrices.

\textbf{19.} $\rho$ toepassen: $\rho(\varepsilon(t)) =
\rho(s(R(t)))\,\rho(q(t))^{-1} = R(t)R(t)^{-1} =
\operatorname{id}$, dus $\varepsilon(t) \in \ker\rho = \{\pm1\}$
(vraag 4). Als product van de \omterm{def:b3:topology:continuity}{continue afbeeldingen} $t \mapsto
s(R(t))$ en $t \mapsto q(t)^{-1} = \bar{q(t)}$ is $\varepsilon$
\omterm{def:b3:topology:continuity}{continu} op het \omterm{def:b3:topology:connected}{samenhangende} interval $\intcc0{2\pi}$ met waarden
in het discrete paar $\{\pm1\}$: ze is dus constant, zeg
$\varepsilon(t) \equiv \varepsilon$. Maar $R(0) = R(2\pi) = I$,
dus $s(R(0)) = s(R(2\pi))$, terwijl $s(R(0)) = \varepsilon\,q(0)
= \varepsilon$ en $s(R(2\pi)) = \varepsilon\,q(2\pi) =
-\varepsilon$: tegenspraak. Er bestaat geen \omterm{def:b3:topology:continuity}{continue} sectie: de
tekenambiguïteit $\pm q$ is globaal, geen gebrek van een
bepaalde formule.

\textbf{20.} \emph{Surjectief:} elke $R \in SO(3)$ is $\eu^{\theta
A_u}$ voor een eenheid $u$ en een $\theta \in \intcc0{2\pi}$
(vragen 6 en 8); is $\theta > \pi$, dan geeft Rodrigues
$\eu^{\theta A_u} = \eu^{(2\pi - \theta)A_{-u}}$ (beide zijn $I +
\sin\theta A_u + (1 - \cos\theta)A_u^2$, en $A_{-u} = -A_u$ met
$\sin(2\pi - \theta) = -\sin\theta$ en $\cos(2\pi - \theta) =
\cos\theta$), dus $R = E(v)$ met $\norm v \leq \pi$.
\emph{Injectief binnenin:} is $E(v) = E(v') \neq I$ met $\norm v,
\norm{v'} < \pi$, dan wint vraag 13 uit het spoor dezelfde hoek
$\theta = \norm v = \norm{v'} \in \intoo0\pi$ terug en, omdat
$\sin\theta \neq 0$, uit het antisymmetrische deel dezelfde as:
$v = v'$; en $E(v) = I$ dwingt op de open bal $\theta \in \{0\}$
af. \emph{Rand:} $E(\pi u) = I + 2A_u^2 = 2uu^{\mathsf T} - I$
hangt van $u$ enkel af via $uu^{\mathsf T}$, waaruit $E(\pi u) =
E(-\pi u)$; omgekeerd geeft $2uu^{\mathsf T} - I = 2u'u'^{\mathsf
T} - I$ toegepast op $u$ dat $u = \langle u', u\rangle u'$, dus
$u' = \pm u$. Inwendig en rand botsen nooit (spoor $> -1$
tegenover $= -1$). Dus induceert $E$ een \omterm{def:b3:topology:continuity}{continue} bijectie van de
bal met de antipodale verlijming --- \omterm{def:b3:topology:compact}{compact} --- op $SO(3)$: een
\omterm{def:b3:topology:continuity}{homeomorfisme}, en $SO(3) \cong \mathbb{RP}^3$. Een middellijn van
$\pi u$ naar $-\pi u$ heeft gelijmde uiteinden: het is een lus in
$SO(3)$, en haar beschrijving met $E$ komt overeen met de familie
rotaties om $u$ uit vraag 10 die een volledige omwenteling
doorloopt.

\textbf{21.} $\rho$ is een morfisme en $\rho_p^{-1} =
\rho_{p^{-1}} = \rho_{\bar p}$, dus $\rho_p\rho_q\rho_p^{-1} =
\rho_{pq\bar p}$; volgens vraag 16 behoudt het vervoegen van een
rotatie haar hoek en draait het haar as met $\rho_p$. Involuties:
$\rho_q^2 = \rho_{q^2} = \operatorname{id}$ dan en slechts dan
als $q^2 = \pm1$. Is $q^2 = 1$, dan is $(q - 1)(q + 1) = q^2 - 1
= 0$ (centrale scalairen, dus die ontbinding is geldig) en $q =
\pm1$ in het delingslichaam $\mathbb H$, wat $\rho_q = I$ geeft:
uitgesloten; $q^2 = -1$ met $q = a + su$ geeft $a^2 - s^2 +
2as\,u = -1$, dus $a = 0$ en $s = 1$: $q$ is een pure eenheid $w$
en $\rho_w$ is de halve slag om $w$ (hoek $\pi$, vraag 5).
\omterm{ex:b3:groups:actions}{Centrum}: commuteert $\rho_q$ met elke $\rho_p$, dan is
$\rho_{pq\bar p} = \rho_q$, dus $pq\bar p = \varepsilon(p)\,q$
met $\varepsilon(p) \in \{\pm1\}$; en $p \mapsto \varepsilon(p) =
(pq\bar p)q^{-1}$ is \omterm{def:b3:topology:continuity}{continu} op het \omterm{def:b3:topology:connected}{samenhangende} $S^3$ en gelijk
aan $1$ in $p = 1$, dus $\equiv 1$: $q$ commuteert met heel
$S^3$, en dus met heel $\mathbb H$ (herschaal), zodat $q \in \R
\cap S^3 = \{\pm1\}$ (vraag 1) en $\rho_q = I$: het \omterm{ex:b3:groups:actions}{centrum} is
triviaal.

\textbf{22.} Omdat $u \perp w$ pure eenheden zijn, is $uw =
u\wedge w$ puur (vraag 18), dus is
\[
w' = qw = \cos\tfrac\theta2\,w + \sin\tfrac\theta2\,u\wedge
w
\]
puur, met norm $\abs q\abs w = 1$. Dan is $w'w = qw\cdot w = qw^2
= -q$, en
\[
\rho_{w'}\rho_w = \rho_{w'w} = \rho_{-q} = \rho_q .
\]
Beide assen $w$ en $w' = \cos\frac\theta2 w +
\sin\frac\theta2(u\wedge w)$ liggen in het vlak $u^\perp$ dat
loodrecht op de rotatieas staat, en $\langle w', w\rangle =
\cos\frac\theta2$: ze maken de halve hoek $\frac\theta2$. Dat is
de klassieke voortbrenging: twee halve slagen om assen die een
hoek $\frac\theta2$ maken, stellen samen tot de rotatie over hoek
$\theta$ om hun gemeenschappelijke loodlijn.

\textbf{23.} De \omterm{def:b3:topology:compact}{compactheid} is vraag 7. Wegsamenhang: $SO(3) =
\rho(S^3)$ is het \omterm{def:b3:topology:continuity}{continue} beeld van de \omterm{def:b3:topology:pathconnected}{wegsamenhangende} sfeer;
of ook: voor $R = \eu^{A}$ met $A$ antisymmetrisch (vraag 8) is
$t \mapsto \eu^{tA}$ een weg in $SO(3)$ van $I$ naar $R$
(orthogonaal want $(\eu^{tA})^{\mathsf T} = \eu^{-tA}$, met
determinant $1$ wegens de \omterm{def:b3:topology:continuity}{continuïteit} vanaf $t = 0$). Voor
$O(3)$: $\det$ is \omterm{def:b3:topology:continuity}{continu} op $\{\pm1\}$, dus is $O(3)$ niet
\omterm{def:b3:topology:connected}{samenhangend}, $O(3) = SO(3) \sqcup D\,SO(3)$ voor een vaste $D$
met $\det D = -1$ (bijvoorbeeld $D = -I$), en
linksvermenigvuldigen met $D$ is een \omterm{def:b3:topology:continuity}{homeomorfisme}: precies twee
componenten, elk een kopie van $SO(3)$. De twee-op-één-afbeelding
$\rho$, zonder sectie volgens vraag 19, is dus een eerlijke
dubbele overdekking van een \omterm{def:b3:topology:connected}{samenhangende} \omterm{def:b3:topology:compact}{compacte} groep door het
\omterm{def:b3:conformal:simplyconnected}{enkelvoudig samenhangende} $S^3$ --- de meetkunde achter de
riemtruc.

\textbf{24.} (i) Matrices:
\[
R_1 = \begin{pmatrix} 0 & -1 & 0\\ 1 & 0 & 0\\ 0 & 0 & 1
\end{pmatrix},
\quad
R_2 = \begin{pmatrix} 1 & 0 & 0\\ 0 & 0 & -1\\ 0 & 1 & 0
\end{pmatrix},
\quad
R_2R_1 = \begin{pmatrix} 0 & -1 & 0\\ 0 & 0 & -1\\ 1 & 0 &
0\end{pmatrix}.
\]
Uit spoor $0 = 1 + 2\cos\theta$ volgt $\cos\theta = -\frac12$:
$\theta = \frac{2\pi}3$. Het antisymmetrische deel $\frac{R -
R^{\mathsf T}}2$ codeert de as als $\sin\theta\,(v_3, -v_2,
v_1)$: hier is $\frac{R - R^{\mathsf T}}2 =
\frac12\bigl(\begin{smallmatrix}0 & -1 & -1\\ 1 & 0 & -1\\ 1
& 1 & 0\end{smallmatrix}\bigr)$, wat (met het woordenboek $A_v$
uit Deel V) $v\sin\theta = \frac12(1, -1, 1)$ leest; en met
$\sin\frac{2\pi}3 = \frac{\sqrt3}2$: $v = \frac1{\sqrt3}(1, -1,
1)$.
(ii) Quaternionen: $q_1 = \cos\frac\pi4 + \sin\frac\pi4\,k =
\frac{\sqrt2}2(1 + k)$, $q_2 = \frac{\sqrt2}2(1 + i)$, en
\[
q_2q_1 = \tfrac12(1 + i)(1 + k) = \tfrac12(1 + k + i + ik)
= \tfrac12\bigl(1 + i - j + k\bigr)
\]
(want $ik = -j$). Dus $\cos\frac\theta2 = \frac12$: $\theta =
\frac{2\pi}3$, en het vectordeel $\frac12(i - j + k)$ heeft
richting $\frac1{\sqrt3}(1, -1, 1)$ --- hetzelfde antwoord,
waarbij de quaternionenweg één regel vermenigvuldiging vergt in
plaats van een matrixproduct: de praktische reden waarom
vluchtsoftware standen in $S^3$ samenstelt.

\textbf{25.} $I + K$ inverteerbaar: uit $(I + K)v = 0$ volgt $0 =
\langle v, v\rangle + \langle v, Kv\rangle = \norm v^2$ (de
antisymmetrie doodt de tweede term): dus $v = 0$. Orthogonaliteit
van $C = C(K)$: met $(I \pm K)^{\mathsf T} = I \mp K$ en het feit
dat alle vier de matrices $I \pm K$ en $(I \pm K)^{-1}$
commuteren (veeltermuitdrukkingen in $K$, plus limieten):
\[
C^{\mathsf T}C = (I + K)^{-\mathsf T}(I - K)^{\mathsf T}
(I - K)(I + K)^{-1}
= (I - K)^{-1}(I + K)(I - K)(I + K)^{-1} = I .
\]
Determinant: $\det(I - K) = \det\bigl((I - K)^{\mathsf
T}\bigr) = \det(I + K)$, dus $\det C = 1$: $C \in SO(n)$. Geen
eigenwaarde $-1$: $Cv = -v$ betekent $(I - K)w = -(I + K)w$ voor
$w = (I + K)^{-1}v$, dat wil zeggen $2w = 0$: dus $v = 0$.
Omkering: uit $C(I + K) = I - K$ los je $K(I + C) = I - C$ op;
omdat $-1 \notin \operatorname{Sp}C$, is $I + C$ inverteerbaar en
$K = (I - C)(I + C)^{-1}$, dat antisymmetrisch is zodra $C$
orthogonaal is zonder eigenwaarde $-1$ (transponeer de
uitdrukking en gebruik $C^{\mathsf T} = C^{-1}$: $K^{\mathsf T} =
(I - C^{-1})(I + C^{-1})^{-1} = (C - I)(C + I)^{-1} = -K$). De
twee afbeeldingen zijn per constructie elkaars inverse: een
globale rationale parametrisering van het \omterm{def:b3:topology:interior}{dichte} open stuk van
$SO(n)$ dat de eigenwaarde $-1$ mijdt --- geen reeksen, geen
goniometrie, en in dimensie $3$ is het de vermomde substitutie
met de halve hoek $K = \tan\frac\theta2\,A_v$.
\end{solution}
